% GENERATED FILE. Edit canonical lessons, then run npm run generate:books.
% canonical-source-hash: fnv1a64:d74fa8eff942cb0a
% canonical-lessons: SA-C157-kupitah, SA-C157-bhitah, SA-C157-cakitah, SA-C157-srantah, SA-C157-ksudhitah

\chapter{Angry, Afraid, Startled, Tired, Hungry}
\label{ch:sa-angry-afraid-startled-tired-hungry}
\hlchaptermodality{157}

\begin{chapteropening}
\textbf{By the end of this chapter:} \emph{I can say angry, afraid, startled, tired and hungry.}

Complete the last lesson of chapter 157: I can say angry, afraid, startled, tired and hungry.
\end{chapteropening}

\section[\texorpdfstring{\sk{कुपितः} (kupitaḥ)}{कुपितः (kupitaḥ)}]{\sk{कुपितः} (kupitaḥ) — angry}

\label{lesson:SA-C157-kupitah}



\begin{quote}
Before the new one: say the Sanskrit for fierce, then the Sanskrit for cruel.
\end{quote}

\subsection*{You'll want to know: \sk{कुपितः}}
\textbf{\sk{कुपितः}} — \emph{kupitaḥ} — \textquotedblleft{}angry\textquotedblright{}.

\begin{grammarlens}[title={What you've built}]
One more word for this chapter.
\end{grammarlens}

\subsection*{Your turn}
\begin{itemize}
\raggedright
  \item \emph{Say it:} \emph{kupitaḥ}
  \item \emph{Say it:} \emph{kupitaḥ}, once more
  \item \emph{Say it:} say \emph{kupitaḥ}
  \item \emph{Recall:} say the Sanskrit for fierce, then the Sanskrit for cruel, then say \emph{kupitaḥ} again
\end{itemize}

\begin{tcolorbox}[breakable,colback=teal!4,colframe=teal!35!black,title={Before you move on}]
What does \sk{कुपितः} mean? (\textquotedblleft{}Angry\textquotedblright{}.) Say it once more.
\end{tcolorbox}
\section[\texorpdfstring{\sk{भीतः} (bhītaḥ)}{भीतः (bhītaḥ)}]{\sk{भीतः} (bhītaḥ) — afraid}

\label{lesson:SA-C157-bhitah}



\begin{quote}
Before the new one: say the Sanskrit for cruel, then the Sanskrit for angry.
\end{quote}

\subsection*{You'll want to know: \sk{भीतः}}
\textbf{\sk{भीतः}} — \emph{bhītaḥ} — \textquotedblleft{}afraid\textquotedblright{}.

\begin{grammarlens}[title={What you've built}]
One more word for this chapter.
\end{grammarlens}

\subsection*{Your turn}
\begin{itemize}
\raggedright
  \item \emph{Say it:} \emph{bhītaḥ}
  \item \emph{Say it:} \emph{bhītaḥ}, once more
  \item \emph{Say it:} say \emph{bhītaḥ}
  \item \emph{Recall:} say the Sanskrit for cruel, then the Sanskrit for angry, then say \emph{bhītaḥ} again
\end{itemize}

\begin{tcolorbox}[breakable,colback=teal!4,colframe=teal!35!black,title={Before you move on}]
What does \sk{भीतः} mean? (\textquotedblleft{}Afraid\textquotedblright{}.) Say it once more.
\end{tcolorbox}
\section[\texorpdfstring{\sk{चकितः} (cakitaḥ)}{चकितः (cakitaḥ)}]{\sk{चकितः} (cakitaḥ) — startled, surprised}

\label{lesson:SA-C157-cakitah}



\begin{quote}
Before the new one: say the Sanskrit for angry, then the Sanskrit for afraid.
\end{quote}

\subsection*{You'll want to know: \sk{चकितः}}
\textbf{\sk{चकितः}} — \emph{cakitaḥ} — \textquotedblleft{}startled, surprised\textquotedblright{}.

\begin{grammarlens}[title={What you've built}]
One more word for this chapter.
\end{grammarlens}

\subsection*{Your turn}
\begin{itemize}
\raggedright
  \item \emph{Say it:} \emph{cakitaḥ}
  \item \emph{Say it:} \emph{cakitaḥ}, once more
  \item \emph{Say it:} say \emph{cakitaḥ}
  \item \emph{Recall:} say the Sanskrit for angry, then the Sanskrit for afraid, then say \emph{cakitaḥ} again
\end{itemize}

\begin{tcolorbox}[breakable,colback=teal!4,colframe=teal!35!black,title={Before you move on}]
What does \sk{चकितः} mean? (\textquotedblleft{}Startled, surprised\textquotedblright{}.) Say it once more.
\end{tcolorbox}
\section[\texorpdfstring{\sk{श्रान्तः} (śrāntaḥ)}{श्रान्तः (śrāntaḥ)}]{\sk{श्रान्तः} (śrāntaḥ) — tired}

\label{lesson:SA-C157-srantah}



\begin{quote}
Before the new one: say the Sanskrit for afraid, then the Sanskrit for startled.
\end{quote}

\subsection*{You'll want to know: \sk{श्रान्तः}}
\textbf{\sk{श्रान्तः}} — \emph{śrāntaḥ} — \textquotedblleft{}tired\textquotedblright{}.

\begin{grammarlens}[title={What you've built}]
One more word for this chapter.
\end{grammarlens}

\subsection*{Your turn}
\begin{itemize}
\raggedright
  \item \emph{Say it:} \emph{śrāntaḥ}
  \item \emph{Say it:} \emph{śrāntaḥ}, once more
  \item \emph{Say it:} say \emph{śrāntaḥ}
  \item \emph{Recall:} say the Sanskrit for afraid, then the Sanskrit for startled, then say \emph{śrāntaḥ} again
\end{itemize}

\begin{tcolorbox}[breakable,colback=teal!4,colframe=teal!35!black,title={Before you move on}]
What does \sk{श्रान्तः} mean? (\textquotedblleft{}Tired\textquotedblright{}.) Say it once more.
\end{tcolorbox}
\section[\texorpdfstring{\sk{क्षुधितः} (kṣudhitaḥ)}{क्षुधितः (kṣudhitaḥ)}]{\sk{क्षुधितः} (kṣudhitaḥ) — hungry}

\label{lesson:SA-C157-ksudhitah}



\begin{quote}
Before the new one: say the Sanskrit for startled, then the Sanskrit for tired.
\end{quote}

\subsection*{You'll want to know: \sk{क्षुधितः}}
\textbf{\sk{क्षुधितः}} — \emph{kṣudhitaḥ} — \textquotedblleft{}hungry\textquotedblright{}.

\begin{grammarlens}[title={What you've built}]
One more word for this chapter.
\end{grammarlens}

\subsection*{Your turn}
\begin{itemize}
\raggedright
  \item \emph{Say it:} \emph{kṣudhitaḥ}
  \item \emph{Say it:} \emph{kṣudhitaḥ}, once more
  \item \emph{Say it:} say \emph{kṣudhitaḥ}
  \item \emph{Recall:} say the Sanskrit for startled, then the Sanskrit for tired, then say \emph{kṣudhitaḥ} again
\end{itemize}

\begin{tcolorbox}[breakable,colback=teal!4,colframe=teal!35!black,title={Before you move on}]
What does \sk{क्षुधितः} mean? (\textquotedblleft{}Hungry\textquotedblright{}.) Say it once more.
\end{tcolorbox}
